% Relatório de diagnóstico do acervo de uma biblioteca
% Modelo do ÉguaTex (https://www.eguatex.com.br), em domínio público pela licença CC0 1.0:
% use, modifique e distribua sem pedir permissão e sem citar a origem.
%
% Compile com pdfLaTeX e BibTeX (o ÉguaTex faz isso sozinho).
% Arquivos: main.tex (este texto), acervo.dat (contagem por classe) e referencias.bib.

\documentclass[11pt,a4paper]{report}

% ---------- idioma, fontes e página ----------
\usepackage[T1]{fontenc}
\usepackage[utf8]{inputenc}
\usepackage[brazilian]{babel}
\usepackage{tgheros}
\renewcommand{\familydefault}{\sfdefault}   % texto todo em fonte sem serifa
\usepackage{microtype}
\usepackage[a4paper,margin=2.5cm]{geometry}
\usepackage{setspace}
\setstretch{1.25}
\setlength{\parindent}{0pt}
\setlength{\parskip}{0.6em}

% ---------- tabelas, gráficos e cores ----------
\usepackage[table]{xcolor}
\usepackage{booktabs}
\usepackage{caption}
\captionsetup{font=small, labelfont={bf,color=azul}, labelsep=period}
\usepackage{tikz}
\usepackage{pgf-pie}
\usepackage{pgfplots}
\usepackage{pgfplotstable}
\pgfplotsset{compat=1.18}

\definecolor{azul}{RGB}{31,78,121}
\definecolor{bom}{RGB}{76,140,74}
\definecolor{regular}{RGB}{230,170,50}
\definecolor{ruim}{RGB}{192,57,43}
\definecolor{linha}{RGB}{236,242,248}

% ---------- títulos das seções em azul, com filete ----------
\renewcommand{\thesection}{\arabic{section}}
\usepackage{titlesec}
\titleformat{\section}{\Large\bfseries\color{azul}}{\thesection}{0.7em}{}[{\color{azul}\titlerule[0.8pt]}]
\titleformat{\subsection}{\large\bfseries\color{azul}}{\thesubsection}{0.6em}{}

% ---------- links e referências no padrão ABNT ----------
% O hyperref vem antes do abntex2cite, senão as citações quebram.
\usepackage[hidelinks]{hyperref}
\usepackage[alf]{abntex2cite}
\addto\captionsbrazilian{\renewcommand{\bibname}{Referências}}

% ---------- instruções de preenchimento ----------
% Os textos em cinza explicam o que escrever em cada parte.
% Troque \instrucoestrue por \instrucoesfalse para esconder todos de uma vez.
\newif\ifinstrucoes
\instrucoestrue
\newcommand{\instrucao}[1]{\ifinstrucoes{\color{gray}\small\itshape #1\par}\fi}

% ---------- dados da capa: preencha aqui ----------
\newcommand{\universidade}{Universidade Federal Exemplo}
\newcommand{\unidade}{Biblioteca Setorial do Instituto de Ciências Exatas}
\newcommand{\titulo}{Diagnóstico do acervo}
\newcommand{\subtitulo}{estado de conservação e composição por área do conhecimento}
\newcommand{\autores}{Nome do autor}
\newcommand{\local}{Cidade}
\newcommand{\data}{Ano}

\hypersetup{pdftitle={\titulo}, pdfauthor={\autores}}

% ---------- dados do levantamento (arquivo acervo.dat) ----------
% A tabela soma as colunas sozinha; os gráficos usam os totais abaixo.
% Ao trocar os dados, atualize também estes totais.
\pgfplotstableread{acervo.dat}\acervo
\pgfplotstablecreatecol[create col/expr={\thisrow{bom}+\thisrow{regular}+\thisrow{ruim}}]{total}\acervo
\pgfplotstablecreatecol[create col/expr={100*\thisrow{ruim}/(\thisrow{bom}+\thisrow{regular}+\thisrow{ruim})}]{pruim}\acervo
\newcommand{\totalbom}{2800}
\newcommand{\totalregular}{870}
\newcommand{\totalruim}{330}
\newcommand{\totalacervo}{4000}

\begin{document}

% ======================================================================
% Capa
% ======================================================================
\begin{titlepage}
  \begin{tikzpicture}[remember picture, overlay]
    \fill[azul] (current page.north west) rectangle ([yshift=-10.5cm]current page.north east);
  \end{tikzpicture}
  \color{white}
  {\large\bfseries\MakeUppercase{\universidade}\par}
  {\unidade\par}
  \vspace{2.2cm}
  {\fontsize{30}{34}\selectfont\bfseries\titulo\par}
  \vspace{0.4em}
  {\Large\subtitulo\par}
  \color{black}
  \vspace{3cm}
  \begin{center}
    \begin{tikzpicture}[scale=0.9]
      % Prateleira estilizada: lombadas nas cores do estado de conservação
      \foreach \x/\c/\h in {0/bom/2.2, 0.55/bom/2.6, 1.1/regular/2.0, 1.6/bom/2.4,
                            2.15/ruim/1.8, 2.6/bom/2.5, 3.15/regular/2.2, 3.7/bom/2.7,
                            4.25/bom/2.1, 4.8/regular/2.4, 5.35/bom/2.6}
        \fill[\c, rounded corners=1pt] (\x,0) rectangle ++(0.45,\h);
      \fill[azul] (-0.3,-0.25) rectangle (6.1,0);
    \end{tikzpicture}
  \end{center}
  \vfill
  {\large\autores\par}
  \vspace{1em}
  {\local, \data\par}
\end{titlepage}

\tableofcontents
\clearpage

% ======================================================================
\section{Apresentação}
% ======================================================================
\instrucao{Diga em um parágrafo qual biblioteca foi avaliada, quem pediu o diagnóstico e para que ele vai servir. Exemplo: subsidiar a política de desenvolvimento de coleções.}

Este relatório apresenta o diagnóstico do acervo da \unidade, com a composição da coleção por área do conhecimento e o estado de conservação dos itens. O resultado orienta as decisões de aquisição, restauração e descarte, de modo que cada leitor encontre o seu livro, como propõe a segunda lei de \citeonline{ranganathan2009}.

% ======================================================================
\section{Metodologia}
% ======================================================================
\instrucao{Explique como os itens foram contados e avaliados: censo ou amostra, período, critérios de cada categoria e quem fez a avaliação. Uma tabela de critérios ajuda quem for repetir o levantamento.}

O levantamento contou todos os itens da coleção de empréstimo, agrupados pelas classes principais da Classificação Decimal de Dewey (CDD). Cada item recebeu uma das três categorias da Tabela~\ref{tab:criterios}.

\begin{table}[htbp]
  \centering
  \caption{Critérios usados para classificar o estado de conservação.}
  \label{tab:criterios}
  \renewcommand{\arraystretch}{1.3}
  \begin{tabular}{>{\bfseries}l p{10.5cm}}
    \toprule
    Estado & Critério \\
    \midrule
    {\color{bom}Bom}         & Capa, lombada e páginas íntegras; pode circular sem restrição. \\
    {\color{regular}Regular} & Desgaste na capa ou na lombada, páginas soltas ou anotações; precisa de pequeno reparo. \\
    {\color{ruim}Ruim}       & Páginas faltando, fungos, umidade ou encadernação perdida; indicar para restauração ou descarte. \\
    \bottomrule
  \end{tabular}
\end{table}

% ======================================================================
\section{Resultados}
% ======================================================================
\instrucao{Mostre os números em tabela e resuma as proporções em gráficos. Comente as classes com maior proporção de itens em estado ruim.}

\subsection{Estado do acervo por classe}

A Tabela~\ref{tab:estado} reúne a contagem por classe da CDD. Foram avaliados \totalacervo{} itens.

\begin{table}[htbp]
  \centering
  \caption{Estado de conservação por classe da CDD (dados de exemplo).}
  \label{tab:estado}
  \rowcolors{2}{linha}{white}
  \pgfplotstabletypeset[
    columns={classe,area,bom,regular,ruim,total,pruim},
    string type,
    assume math mode=true, % números na fonte do texto, sem modo matemático
    every head row/.style={before row=\toprule, after row=\midrule},
    every last row/.style={after row=\bottomrule},
    columns/classe/.style={column name={CDD}, column type=l},
    columns/area/.style={column name={Área}, column type=l},
    columns/bom/.style={column name={Bom}, column type=r},
    columns/regular/.style={column name={Regular}, column type=r},
    columns/ruim/.style={column name={Ruim}, column type=r},
    columns/total/.style={column name={Total}, column type=r, numeric type, fixed, precision=0, 1000 sep={.}},
    columns/pruim/.style={column name={\% ruim}, column type=r, numeric type, fixed, fixed zerofill, precision=1, use comma},
  ]\acervo
\end{table}

\subsection{Composição e conservação}

A Figura~\ref{fig:pizzas} resume a coleção em duas proporções: a parte de cada grande área no total e a distribuição geral do estado de conservação.

\begin{figure}[htbp]
  \centering
  \begin{tikzpicture}
    % Composição por área (contagens; sum=auto converte em setores proporcionais)
    \pie[sum=auto, radius=1.55, text=legend, after number=,
         color={azul!90, azul!65, azul!45, azul!28, gray!25}]
      {900/Literatura (800), 760/Ciências naturais (500), 640/Ciências sociais (300),
       540/Tecnologia (600), 1160/Outras classes}
  \end{tikzpicture}
  \hfill
  \begin{tikzpicture}
    % Estado de conservação (contagens; o percentual vai na legenda)
    \pie[sum=auto, radius=1.55, text=legend, after number=, color={bom, regular, ruim}]
      {2800/Bom (70\,\%), 870/Regular (21{,}75\,\%), 330/Ruim (8{,}25\,\%)}
  \end{tikzpicture}
  \caption{À esquerda, número de itens por área; à direita, número de itens por estado de conservação.}
  \label{fig:pizzas}
\end{figure}

\instrucao{Destaque a classe mais crítica com o número e o percentual, como no exemplo abaixo.}
Dos \totalacervo{} itens, \totalbom{} estão em bom estado, \totalregular{} em estado regular e \totalruim{} em estado ruim. A classe 500 tem a maior proporção de itens em estado ruim, o que se explica pelo uso intenso dos livros-texto de Física e Matemática.

% ======================================================================
\section{Recomendações}
% ======================================================================
\instrucao{Transforme os resultados em ações, com prazo e responsável. Uma lista numerada por prioridade funciona bem.}

\begin{enumerate}
  \item Encaminhar os itens em estado ruim da classe 500 para avaliação de restauração ou reposição.
  \item Fazer pequenos reparos nos itens em estado regular durante o período de férias.
  \item Revisar a política de desenvolvimento de coleções com base na composição por área.
\end{enumerate}

% ======================================================================
% Referências
% ======================================================================
\begingroup
\let\chapter\section
\phantomsection
\addcontentsline{toc}{section}{\bibname}
\bibliography{referencias}
\endgroup

\end{document}
